\documentclass[11pt]{article}
\usepackage[margin=1in]{geometry}
\usepackage{amsmath,amssymb,amsthm,booktabs,tabularx}
\usepackage{microtype}
\newtheorem{proposition}{Proposition}
\newtheorem{lemma}{Lemma}
\newtheorem{corollary}{Corollary}
\newtheorem{definition}{Definition}
\newcommand{\dd}{\mathrm d}

\begin{document}

\section*{Constrained efficiency and infinite-horizon dynamics}

Two questions remain after solving the household problems.  First, does the
marginal-value gap imply inefficiency under the institutions households
actually face?  Second, can the finite transition be replaced by an exact
infinite-horizon result?  The answers are different.  The baseline model gives
a sharp but limited efficiency result.  An infinite-horizon theorem is
available for a smooth specialization, but the realization-tax kink prevents
that theorem from applying directly to the model as written.

\subsection*{1. Efficiency under the baseline institutions}

Fix a date and hold fertility, cohort masses, tenure, prices, and the identity
of future households fixed.  A \emph{baseline-implementable} perturbation must
respect three institutional restrictions.  A buyer cannot borrow against
current income or the rebate to meet the down payment; the common rebate
arrives after closing; and housing retained by an incumbent cannot be occupied
by another household.  The government may change only the common transfers
implied by its period-by-period budget.  In particular, it cannot make a
buyer-specific transfer at closing.  Let $\mathcal F_t^B$ denote the set of
allocations attainable with these instruments, the household mortgage and bond
constraints, and the aggregate goods and housing constraints.  A constrained
planner may select any allocation in $\mathcal F_t^B$, but may not reassign
pre-closing liquid wealth.  An allocation is locally constrained Pareto
efficient if no differentiable path in $\mathcal F_t^B$ makes every household
alive at $t$ weakly better off over its remaining lifetime and at least one
strictly better off.

For a young owner, define liquid wealth left at closing by
\begin{equation}
 L_{it}=b_i-(1-\phi)P_th_{it}^{O}\geq 0.
 \label{eq:liquidity_slack}
\end{equation}
The multiplier on this inequality, expressed in current consumption goods, is
the financing component of the access wedge,
\begin{equation}
 \zeta_{it}^{O,F}=\frac{\mu_{it}}{\lambda_{it}}(1-\phi)P_t.
 \label{eq:financing_wedge}
\end{equation}
It measures the value of additional liquidity at closing, not the gain from a
reallocation that is feasible under the original credit technology.

\begin{lemma}[The constrained buyer is not a feasible recipient]
\label{lem:tangent}
Suppose \eqref{eq:liquidity_slack} binds.  Along any one-sided differentiable
baseline-implementable perturbation that holds $P_t$ and pre-closing liquid
wealth $b_i$ fixed, the buyer's housing change satisfies
$\dot h_{it}^{O}\leq0$.
\end{lemma}

\begin{proof}
Feasibility requires $L_{it}(\epsilon)\geq0$ for $\epsilon\geq0$ and
$L_{it}(0)=0$.  Hence
\[
 0\leq \dot L_{it}=-(1-\phi)P_t\dot h_{it}^{O},
\]
which gives the result.  A change in the common date-$t$ rebate does not enter
this derivative because the rebate is unavailable at closing.
\end{proof}

The lemma changes the interpretation of the steady-state value gap.  At a
positive steady state, realization taxes are zero and the comparison between a
constrained young owner and an interior old incumbent is
\begin{equation}
 MV_{it}^{Y}-MV_{jt}^{O}=\zeta_{it}^{O,F}>0.
 \label{eq:ss_gap}
\end{equation}
But the perturbation that increases $h_{it}^{O}$ is outside the tangent set of
baseline-implementable allocations.  The positive multiplier is therefore
consistent with the first-order conditions of a planner constrained by the
same closing rule.  Equation~\eqref{eq:ss_gap} does not prove that the
competitive allocation is constrained inefficient.

This conclusion is specific to the financing component of the gap.  During an
appreciation episode, the realization tax can distort a transfer that is
feasible without relaxing credit.

\begin{proposition}[A baseline transition inefficiency]
\label{prop:baseline_tax}
Suppose households form an atomless population and $\ell_t>0$.  Suppose there
are matched subsets of young owners and old incumbents, each of mass $m>0$, on
which allocations lie in a compact interior set, the relevant objectives are
twice continuously differentiable, and the following slackness margins are
bounded away from zero.  Young saving, the
down-payment constraint, the owner-size constraint, and the young buyer's
future old-age retention and estate constraints are slack; old incumbents have
interior retention and estate choices.  Suppose small title transfers use no
real resources.  Then the date-$t$ competitive allocation is not locally
Pareto efficient in $\mathcal F_t^B$ when fertility and population are fixed.
\end{proposition}

\begin{proof}
Transfer $\epsilon$ units of housing within each matched pair at the market
price.  The young can finance the purchase because the closing constraint is
slack.  When the buyer becomes old, hold baseline retained housing fixed and
sell the added unit.  Slack future constraints make this path feasible.  The
household first-order conditions imply that both parties' private welfare
changes are zero to first order: the incumbent receives the after-tax value of
the marginal unit, and the unconstrained buyer pays its private user cost.  The
sale raises current gains-tax revenue by $m\ell_t\epsilon$.  The baseline
budget therefore raises the common current transfer by
\[
 \Delta T_t=\frac{m\ell_t\epsilon}{Y_t+O_t}>0.
\]
Every household alive at $t$ can consume this transfer and is therefore better
off unless its private allocation was perturbed.  On the matched subsets,
twice continuous differentiability and uniform slackness make the private
losses $O(\epsilon^2)$ uniformly, whereas the transfer gain is of order
$\epsilon$.  For sufficiently small $\epsilon$, every current household is
weakly better off and some are strictly better off.
\end{proof}

The proposition requires current appreciation.  If title transfer uses
$K_{ijt}$ current goods per unit, let $\bar K_t$ be its average over the matched
pairs.  The common-current-rebate construction then requires
$\ell_t>\bar K_t$.  A larger present-value surplus that arrives only at $t+1$
does not complete this particular construction under the baseline
period-by-period government budget, because the old incumbent cannot be
compensated after death.

\begin{corollary}[No unconditional efficiency classification]
The maintained model class contains both locally constrained-efficient and
locally constrained-inefficient allocations.
\end{corollary}

\begin{proof}
Proposition~\ref{prop:baseline_tax} supplies the inefficient case.  For the
efficient case, consider a stationary allocation with no realization taxes,
allocation-independent tax revenue, and slack housing and financing limits.
If a feasible perturbation made all current households weakly better off and
one strictly better off, household optimality at the common goods, bond, and
housing prices would make the value of their aggregate net purchases strictly
positive.  Market clearing makes the same sum zero, a contradiction.
\end{proof}

Thus an appreciation episode can fail constrained efficiency when an
unconstrained buyer can absorb housing released by an incumbent.  If all
potential recipients are against their closing constraints, the
marginal-value comparison alone is silent.

The smallest additional primitive that makes the steady-state access wedge
actionable is a pledgeable down-payment advance.  Let a lender or the government
advance $(1-\phi)P_t\epsilon$ at closing and recover its market-value repayment
from the buyer's next-period resources or resale proceeds.  This changes the
tangent condition in Lemma~\ref{lem:tangent} without changing present-value
resources.  If
\begin{equation}
 \zeta_{it}^{O,F}>K_{ijt},
 \label{eq:credit_test}
\end{equation}
then a small transfer to the constrained buyer generates enough surplus to
compensate the incumbent and pay the real transfer cost.  This is a theorem
about the value of adding a liquidity instrument, not a theorem that the
baseline competitive allocation violates its own constraints.

\subsection*{2. The autonomous transition system}

The finite-dimensional numerical specialization has seven predetermined
states.  A convenient ordering is
\begin{equation}
 s_t=(Y_t,O_t,\pi_{t-1}^{O},a_t^{R},a_t^{O},H_t^{O},P_{t-1})\in\mathbb R^7.
 \label{eq:states}
\end{equation}
The two nonpredetermined coordinates are the current price and an auxiliary
copy of next period's price, $j_t=(P_t,Q_t)$.  Provided its derivative with
respect to $T_t$ is nonzero, the fiscal equation locally determines the current
transfer.  After substituting this transfer, the seven
state laws, housing clearing, and the identity $Q_t=P_{t+1}$ form nine
autonomous equations:
\begin{equation}
 \mathcal F(x_t,x_{t+1};\vartheta)=0,
 \qquad x_t=(s_t,j_t)\in\mathbb R^9.
 \label{eq:autonomous_system}
\end{equation}
This system is exact for the homogeneous-type specialization used in the
numerical construction.  With continuously distributed or otherwise
unrestricted heterogeneity, the old-state distribution is itself a state and
the corresponding system is generally infinite-dimensional.

The following result states what the usual saddle-path calculation would prove
if the equilibrium map were smooth.

\begin{proposition}[Local infinite-horizon transition in the smooth case]
\label{prop:smooth_transition}
Fix the post-shock value of $\vartheta$ and let $x^*$ be its steady state.
Suppose $\mathcal F$ is continuously differentiable near $(x^*,x^*)$ and
$D_2\mathcal F(x^*,x^*)$ is nonsingular.  Let
\[
 A=-[D_2\mathcal F(x^*,x^*)]^{-1}D_1\mathcal F(x^*,x^*).
\]
Suppose $A$ has seven eigenvalues strictly inside the unit circle and two
strictly outside it.  Finally, suppose the projection of the seven-dimensional
stable eigenspace onto the predetermined coordinates $s$ is nonsingular.  Then
every admissible initial state $s_0$ sufficiently close to $s^*$ selects a unique
$j_0$ such that the resulting equilibrium remains near $x^*$ and converges to
$x^*$.  The transition is locally unique among paths that remain in that
neighborhood.
\end{proposition}

\begin{proof}
The implicit-function theorem applied to
\eqref{eq:autonomous_system} gives a local map $x_{t+1}=g(x_t)$ with
$Dg(x^*)=A$.  Hyperbolicity gives a seven-dimensional local stable manifold.
The nonsingular projection makes this manifold a graph $j=\Gamma(s)$ over the
predetermined states.  Hence $j_0=\Gamma(s_0)$ is the unique local choice whose
orbit converges to $x^*$.
\end{proof}

For a small permanent shock, the state inherited from the old steady state is
close to the new steady state.  Proposition~\ref{prop:smooth_transition} would
therefore turn the local saddle-path conditions into an exact transition
between the two steady states.  It neither imposes a terminal date nor assumes
that terminally closed finite systems converge.

\subsection*{3. Why the proposition does not yet apply}

The statutory gains rule,
\begin{equation}
 \ell_t=\tau^g\max\{P_t-P_{t-1},0\},
 \label{eq:kink}
\end{equation}
is not differentiable at either steady state.  Consequently there is no single
matrix $A$ to which the smooth stable-manifold theorem can be applied.  The
existing diagnostic supplies strong numerical evidence: after eliminating the
static transfer, the next-date Jacobian has rank nine; each of the four
one-sided derivative selections has seven stable and two unstable roots; and
the stable subspace projects with full rank onto the seven predetermined
states.  At the final steady state, the largest stable modulus across the
reported generalized-derivative grid is $0.547$, the smallest unstable modulus
is $3.746$, and the smallest reported projection singular value is $0.544$.

These calculations do not prove an infinite-horizon result.  They use finite
differences in the homogeneous specialization and sample, rather than exhaust,
the generalized derivatives.  More importantly, hyperbolicity of every frozen
matrix does not imply stability under an arbitrary sequence of switches among
those matrices.

For the exact kinked model, a sufficient theorem would require strong
regularity of an infinite-horizon boundary-value operator.  Place deviations
from the final steady state in the weighted sequence space
\begin{equation}
 \mathbb X_\rho=\left\{(x_t)_{t\geq0}:
 \sup_{t\geq0}\rho^{-t}\lVert x_t-x^*\rVert<\infty\right\},
 \qquad \alpha<\rho<1,
 \label{eq:weighted_space}
\end{equation}
where $\alpha$ is a uniform bound on the stable decay rate.  Substitute the
positive-part rule \eqref{eq:kink} directly into the one-date residual and call
the resulting locally Lipschitz map $\mathcal F^{\max}$.  If $\Pi_s$ selects
the seven predetermined coordinates, the relevant operator is
\begin{equation}
 \mathfrak H(\boldsymbol x;s_0,\varepsilon)=
 \left(
 (\mathcal F^{\max}(x_t,x_{t+1};\varepsilon))_{t\geq0},
 \Pi_sx_0-s_0
 \right).
 \label{eq:augmented_operator}
\end{equation}
The initial-state equation is essential.  Without it, the linearized residual
has the seven-dimensional stable manifold in its kernel and cannot be
invertible.

A sufficient nonsmooth theorem requires the following conditions:
\begin{enumerate}
 \item every admissible generalized next-date Jacobian is nonsingular, with a
 common bound on its inverse;
 \item the resulting switched linear system has a uniform exponential
 dichotomy with a seven-dimensional stable bundle and a two-dimensional
 unstable bundle, and stable rate at most $\alpha$; and
 \item the stable bundle projects uniformly and nonsingularly onto the seven
 predetermined states.
\end{enumerate}
Together with a uniform semismooth remainder bound, these conditions make each
admissible generalized derivative of the augmented operator
\eqref{eq:augmented_operator} a bounded bijection, with a common bound on its
inverse.  This is the strong-regularity condition that a nonsmooth implicit-
function or graph-transform argument needs.  It would deliver a unique local
path in $\mathbb X_\rho$, and membership in that space would establish
convergence.  This does not assume convergence of finite-horizon solutions; it
derives convergence from bounded invertibility of the infinite linearized
boundary-value problem.

The existing root-count calculation does not verify this uniform switched-
system condition.  Establishing it would require either an analytic common-cone
argument or a certified bound over all admissible generalized Jacobians and
their products.  For the model with a general distribution of household types,
the same argument would also have to be carried out on a space of measures.

The clean alternative is to make the gains schedule continuously
differentiable at zero gains, including the special case $\tau^g=0$.  Then the
smooth theorem applies once the rank, root-count, and projection conditions are
verified exactly.  That modification removes the lock-in kink that motivates
the transition exercise.  The nonsmooth infinite-horizon proof is therefore
possible in principle, but it is not a minimal addition to the analytical
model.

\subsection*{4. Claims supported by the analysis}

\begin{center}
\small
\begin{tabularx}{\textwidth}{@{}p{0.26\textwidth}X@{}}
\toprule
Status & Claim\\
\midrule
Established & A binding down-payment wedge does not establish inefficiency
relative to the baseline credit technology, because the proposed recipient is
outside the baseline tangent set.\\
Established & Current realization taxation can generate a baseline-feasible
local Pareto improvement when an unconstrained young buyer and an interior old
incumbent can trade and the common current rebate returns the revenue.\\
One additional instrument & A pledgeable down-payment advance turns
$\zeta_{it}^{O,F}>K_{ijt}$ into a local policy test at a steady state.\\
One smoothness change & A differentiable gains schedule plus the stated exact
rank, hyperbolicity, and projection conditions yields a locally determinate
infinite-horizon transition.\\
Not established & The numerical saddle-path diagnostic does not prove strong
regularity of the statutory kinked system or of the model with a distribution
state.\\
\bottomrule
\end{tabularx}
\end{center}

\end{document}
